\documentclass[12pt,a4paper]{article}

\usepackage[T1]{fontenc}
\usepackage[utf8]{inputenc}
\usepackage{lmodern}
\usepackage[ngerman]{babel}

\usepackage{amsmath,amssymb,amsfonts,amsthm}
\usepackage{mathtools}
\usepackage{geometry}
\geometry{margin=2.8cm}

\usepackage{enumitem}
\usepackage{booktabs}
\usepackage{hyperref}

\theoremstyle{plain}
\newtheorem{theorem}{Satz}[section]
\newtheorem{lemma}[theorem]{Lemma}
\newtheorem{corollary}[theorem]{Korollar}
\newtheorem{proposition}[theorem]{Proposition}

\theoremstyle{definition}
\newtheorem{definition}[theorem]{Definition}

\theoremstyle{remark}
\newtheorem{remark}[theorem]{Bemerkung}

\title{Zu einer strukturierten Beweisstrategie für gemeinsame Primteiler von Binomialkoeffizienten}
\author{}
\date{März 2026}

\begin{document}
	
	\maketitle
	
	\begin{abstract}
		Wir untersuchen eine Beweisstrategie für die Aussage, dass für ganze Zahlen
		\[
		1\le i<j\le \frac n2
		\]
		eine Primzahl \(p\ge i\) existiert mit
		\[
		p\mid \gcd\!\left(\binom ni,\binom nj\right).
		\]
		Der vorliegende Text ist als strukturierte klassische Fassung einer Beweisarchitektur formuliert. Der algebraische Teil beruht auf einer Master-Identität für Binomialkoeffizienten, auf der Lokalisierung des tame-Fensters \((j-i,j]\) und auf einem Brückenlemma für hohe Primzahlpotenzen. Der globale Teil reduziert den blockierten Fall für \(i\ge10\) auf die Existenz zweier benachbarter \(S_i\)-glatter Zahlen im \(i\)-Block. Für \(i=3\) wird ein elementares Intervalltrennungsargument angegeben, und für \(4\le i\le9\) wird eine reproduzierbare computationale Zertifizierung formuliert. Der Text markiert ausdrücklich die Stellen, an denen zusätzliche formale Verifikation noch sinnvoll ist.
	\end{abstract}
	
	\tableofcontents
	
	\section{Einleitung}
	
	Wir betrachten die folgende Zielaussage.
	
	\begin{quote}
		Für ganze Zahlen \(1\le i<j\le n/2\) existiert eine Primzahl \(p\ge i\), die sowohl \(\binom ni\) als auch \(\binom nj\) teilt.
	\end{quote}
	
	Die Beweisstrategie zerfällt in mehrere Schichten:
	
	\begin{enumerate}[label=\arabic*)]
		\item Untersuchung des kurzen Blocks
		\[
		P_i(n):=n(n-1)\cdots(n-i+1),
		\]
		\item Separation freier Hochprimzahlen \(>j\),
		\item Lokalisierung der einzigen absorbierbaren Primzahlen im tame-Fenster \((j-i,j]\),
		\item Brückenargument für hohe Primzahlpotenzen,
		\item residuale Kontrolle über ein schwaches Escape-Lemma,
		\item globales Schubfach- und Adjazenzargument für \(i\ge10\),
		\item elementare bzw. endliche Behandlung der kleinen Fälle.
	\end{enumerate}
	
	\section{Die Zielaussage}
	
	\begin{theorem}[Zielaussage]
		\label{thm:target}
		Für ganze Zahlen \(n,i,j\) mit
		\[
		1\le i<j\le \frac n2
		\]
		existiert eine Primzahl \(p\ge i\) mit
		\[
		p\mid \gcd\!\left(\binom ni,\binom nj\right).
		\]
	\end{theorem}
	
	\begin{remark}
		Der vorliegende Text ist als strukturierte Beweisstrategie formuliert. Mehrere Teile sind streng ausgearbeitet, andere werden auf kontrollierte Restbausteine reduziert.
	\end{remark}
	
	\section{Klassische Werkzeuge}
	
	\begin{lemma}[Kummer]
		Für eine Primzahl \(p\) und \(n\ge k\ge0\) ist
		\[
		v_p\!\left(\binom nk\right)
		\]
		gleich der Anzahl der Überträge beim Addieren von \(k\) und \(n-k\) in Basis \(p\).
	\end{lemma}
	
	\begin{lemma}[Legendre]
		Für \(m\ge1\) und eine Primzahl \(p\) gilt
		\[
		v_p(m!)=\sum_{r\ge1}\left\lfloor \frac{m}{p^r}\right\rfloor.
		\]
	\end{lemma}
	
	\begin{lemma}[Sylvester--Schur]
		\label{lem:sylvester-schur}
		Für \(n\ge 2k\) besitzt das Produkt
		\[
		n(n-1)\cdots(n-k+1)
		\]
		einen Primfaktor \(>k\).
	\end{lemma}
	
	\begin{lemma}[Master-Identität]
		\label{lem:master}
		Für \(n\ge j>i\ge1\) gilt
		\[
		\binom nj \binom ji
		=
		\binom ni \binom{n-i}{j-i}.
		\]
		Insbesondere folgt für jede Primzahl \(p\),
		\[
		v_p\!\left(\binom nj\right)
		=
		v_p\!\left(\binom ni\right)
		+
		v_p\!\left(\binom{n-i}{j-i}\right)
		-
		v_p\!\left(\binom ji\right).
		\]
	\end{lemma}
	
	\section{Das tame-Fenster}
	
	\begin{lemma}[Tame-Fenster]
		\label{lem:tame-window}
		Sei \(p>i\) prim. Dann gilt
		\[
		p\mid \binom ji
		\quad\Longleftrightarrow\quad
		p\in(j-i,j].
		\]
	\end{lemma}
	
	\begin{proof}
		Da \(p>i\), tritt \(p\) im Nenner \(i!\) nicht auf. Also teilt \(p\) genau dann \(\binom ji\), wenn das Produkt
		\[
		j(j-1)\cdots(j-i+1)
		\]
		ein Vielfaches von \(p\) enthält. Wegen \(p>i\) kann in einem Block der Länge \(i\) höchstens ein Vielfaches von \(p\) vorkommen. Dies ist genau dann der Fall, wenn \(p\in(j-i,j]\).
	\end{proof}
	
	\begin{lemma}[Einfache tame-Bewertung]
		\label{lem:tame-valuation}
		Ist \(q\in(j-i,j]\) prim und \(q>i\), dann gilt
		\[
		v_q\!\left(\binom ji\right)=1.
		\]
	\end{lemma}
	
	\begin{proof}
		Im Produkt
		\[
		j(j-1)\cdots(j-i+1)
		\]
		liegt wegen \(q>i\) höchstens ein Vielfaches von \(q\), und für \(q\in(j-i,j]\) genau eines. Da \(q\nmid i!\), ist die Bewertung gleich \(1\).
	\end{proof}
	
	\begin{lemma}[Nicht-tame Primteiler liefern Zeugen]
		\label{lem:nontame_witness}
		Sei \(1\le i<j\le n/2\), und sei \(p>i\) eine Primzahl mit
		\[
		p\mid P_i(n):=n(n-1)\cdots(n-i+1).
		\]
		Angenommen,
		\[
		p\notin (j-i,j].
		\]
		Dann ist \(p\) ein Zeuge für \((n,i,j)\), also
		\[
		p\mid \gcd\!\left(\binom ni,\binom nj\right).
		\]
	\end{lemma}
	
	\begin{proof}
		Da \(p>i\), tritt \(p\) im Nenner \(i!\) nicht auf. Aus \(p\mid P_i(n)\) folgt daher
		\[
		p\mid \binom ni.
		\]
		Ferner ist \(p\notin(j-i,j]\). Nach Lemma~\ref{lem:tame-window} folgt also
		\[
		p\nmid \binom ji.
		\]
		Mit Lemma~\ref{lem:master} erhalten wir
		\[
		v_p\!\left(\binom nj\right)
		=
		v_p\!\left(\binom ni\right)
		+
		v_p\!\left(\binom{n-i}{j-i}\right)
		-
		v_p\!\left(\binom ji\right).
		\]
		Da \(v_p\!\left(\binom ni\right)\ge1\) und \(v_p\!\left(\binom ji\right)=0\), folgt
		\[
		v_p\!\left(\binom nj\right)\ge1.
		\]
		Also teilt \(p\) beide Binomialkoeffizienten.
	\end{proof}
	
	\section{Freie Hochmoden und erste Reduktion}
	
	\begin{definition}
		Eine Primzahl \(p\) heißt \emph{Zeuge} für \((n,i,j)\), wenn
		\[
		p\mid \gcd\!\left(\binom ni,\binom nj\right).
		\]
	\end{definition}
	
	\begin{definition}
		Ein Tripel \((n,i,j)\) heißt \emph{blockiert}, wenn kein Zeuge existiert.
	\end{definition}
	
	\begin{lemma}[Freier Hochmodus]
		\label{lem:free-large-prime}
		Besitzt der Block
		\[
		P_i(n)=n(n-1)\cdots(n-i+1)
		\]
		einen Primteiler \(p>j\), so ist \(p\) ein Zeuge.
	\end{lemma}
	
	\begin{proof}
		Da \(p>j\), tritt \(p\) weder in \(i!\) noch in \(j!\) auf. Weil \(p\mid P_i(n)\), folgt
		\[
		p\mid \binom ni.
		\]
		Der \(j\)-Block enthält den \(i\)-Block als Teilprodukt, also ebenfalls den durch \(p\) teilbaren Faktor. Da \(p\nmid j!\), folgt auch
		\[
		p\mid \binom nj.
		\]
	\end{proof}
	
	\begin{corollary}
		\label{cor:j-smooth}
		Ist \((n,i,j)\) blockiert, so sind alle Primteiler von \(P_i(n)\) höchstens \(j\).
	\end{corollary}
	
	\begin{theorem}[Erste Fallzerlegung]
		\label{thm:first-reduction}
		Sei \(1\le i<j\le n/2\). Dann gilt:
		\begin{enumerate}[label=\textup{(\Alph*)}]
			\item Entweder \(P_i(n)\) besitzt einen Primteiler \(p>j\), dann existiert ein Zeuge.
			\item Oder \(P_i(n)\) ist \(j\)-glatt. Dann besitzt \(P_i(n)\) nach Lemma~\ref{lem:sylvester-schur} dennoch einen Primteiler \(q>i\), also einen Primteiler
			\[
			q\in(i,j].
			\]
		\end{enumerate}
	\end{theorem}
	
	\begin{proof}
		Fall (A) ist Lemma~\ref{lem:free-large-prime}. Falls (A) nicht vorliegt, sind alle Primteiler von \(P_i(n)\) höchstens \(j\). Nach Lemma~\ref{lem:sylvester-schur} besitzt \(P_i(n)\) dennoch einen Primteiler \(>i\). Also liegt wenigstens ein Primteiler in \((i,j]\).
	\end{proof}
	
	\section{Das Brückenlemma für hohe Potenzen}
	
	\begin{lemma}[Brückenlemma]
		\label{lem:bridge}
		Sei \(p\) prim mit \(i\le p\le j\). Es gebe ein \(k\in\{0,\dots,i-1\}\) mit
		\[
		p^v\mid (n-k),\qquad v\ge2,\qquad p^v>j.
		\]
		Dann ist \(p\) ein Zeuge für \((n,i,j)\).
	\end{lemma}
	
	\begin{proof}
		Für den \(i\)-Block gilt: Da \(p\ge i\), kann im Block
		\[
		n(n-1)\cdots(n-i+1)
		\]
		höchstens ein Vielfaches von \(p\) auftreten, nämlich \(n-k\). Wegen \(v\ge2\) folgt
		\[
		v_p\!\left(n(n-1)\cdots(n-i+1)\right)\ge 2.
		\]
		Ferner gilt \(v_p(i!)\le1\), da \(p\ge i\). Also
		\[
		v_p\!\left(\binom ni\right)\ge1.
		\]
		
		Für den \(j\)-Block setze
		\[
		N_j:=n(n-1)\cdots(n-j+1).
		\]
		Da \(k\le i-1<j\), liegt der Faktor \(n-k\) auch in \(N_j\). Schreibe \(A_r\) für die Anzahl der Vielfachen von \(p^r\) im \(j\)-Block. Dann gilt
		\[
		v_p(N_j)=\sum_{r\ge1} A_r.
		\]
		Nach Legendre ist
		\[
		v_p(j!)=\sum_{r\ge1}\left\lfloor \frac{j}{p^r}\right\rfloor.
		\]
		
		Aus \(p^v\mid(n-k)\) und \(p^v>j\) folgt, dass der \(j\)-Block genau ein Vielfaches von \(p^v\) enthält. Also
		\[
		A_v=1
		\qquad\text{und}\qquad
		\left\lfloor \frac{j}{p^v}\right\rfloor=0.
		\]
		Ferner gilt für jedes \(r\ge1\),
		\[
		A_r\ge \left\lfloor \frac{j}{p^r}\right\rfloor.
		\]
		Somit
		\[
		v_p(N_j)-v_p(j!)
		=
		\sum_{r\ge1}\left(A_r-\left\lfloor \frac{j}{p^r}\right\rfloor\right)\ge1.
		\]
		Also
		\[
		v_p\!\left(\binom nj\right)\ge1.
		\]
		Damit teilt \(p\) beide Binomialkoeffizienten.
	\end{proof}
	
	\section{Residualfall und schwaches Escape-Lemma}
	
	\begin{definition}[Lonely-Cofaktor]
		Sei \(q\in(j-i,j]\) prim und
		\[
		v_q(n-k_q)=1
		\]
		für ein \(k_q\in\{0,\dots,i-1\}\). Dann schreiben wir
		\[
		n-k_q=q\,m_q,\qquad \gcd(q,m_q)=1.
		\]
		Die Zahl \(m_q\) heißt der \emph{lonely-Cofaktor}.
	\end{definition}
	
	\begin{lemma}[Schwaches Escape-Lemma]
		\label{lem:weak-escape}
		Sei \((n,i,j)\) blockiert, und sei
		\[
		n-k_q=q\,m_q,\qquad q\in(j-i,j],\qquad \gcd(q,m_q)=1,\qquad v_q(n-k_q)=1.
		\]
		Dann gilt: Jeder Primteiler \(r>i\) von \(m_q\) liegt selbst im Intervall
		\[
		r\in(j-i,j].
		\]
		Insbesondere besitzt \(m_q\) keinen Primteiler \(>j\).
	\end{lemma}
	
	\begin{proof}
		Sei \(r\mid m_q\) mit \(r>i\). Dann gilt \(r\mid(n-k_q)\), also tritt \(r\) im \(i\)-Block auf. Da \(r>i\), gilt
		\[
		r\mid \binom ni.
		\]
		
		Wäre \(r>j\), so wäre \(r\) nach Lemma~\ref{lem:free-large-prime} ein Zeuge. Das ist unmöglich, da \((n,i,j)\) blockiert ist. Also gilt \(r\le j\).
		
		Wäre nun \(r\notin(j-i,j]\), so wäre \(r\) nach Lemma~\ref{lem:tame-window} nicht tame. Dann wäre \(r\) nach Lemma~\ref{lem:nontame_witness} ein Zeuge, wiederum ein Widerspruch.
		
		Folglich gilt
		\[
		r\in(j-i,j].
		\]
	\end{proof}
	
	\begin{remark}
		Dieses Lemma ist bewusst schwächer als die Behauptung, \(m_q\) sei vollständig \(i\)-glatt. Für den globalen Schluss genügt, dass alle großen Restprimteiler in dasselbe schmale Fenster gezwungen werden.
	\end{remark}
	
	\section{Der Fall \(i=3\)}
	
	\begin{lemma}[Tame-Fenster für \(i=3\)]
		\label{lem:tame_i3}
		Sei \(j\ge4\). Eine Primzahl \(q>3\) teilt \(\binom j3\) genau dann, wenn
		\[
		q\in(j-3,j].
		\]
		Da \(q\) ganzzahlig und prim ist, bedeutet dies
		\[
		q\in\{j-2,\;j-1,\;j\}\cap\mathbb P.
		\]
	\end{lemma}
	
	\begin{proof}
		Dies ist Lemma~\ref{lem:tame-window} im Spezialfall \(i=3\). Da das Intervall
		\((j-3,j]\) für ganze Zahlen genau die drei Werte \(j-2,j-1,j\) enthält, folgt die zweite Aussage unmittelbar.
	\end{proof}
	
	\begin{lemma}[Kein Tame-Treffer im Dreierblock für \(n>10\)]
		\label{lem:no_tame_hit_i3}
		Seien \(i=3\), \(3<j\le n/2\) und \(n>10\). Dann enthält der Dreierblock
		\[
		\{n,\;n-1,\;n-2\}
		\]
		keine Primzahl aus dem tame-Fenster \((j-3,j]\).
	\end{lemma}
	
	\begin{proof}
		Nach Lemma~\ref{lem:tame_i3} müsste ein solcher Tame-Treffer eine der Zahlen
		\[
		j,\quad j-1,\quad j-2
		\]
		sein.
		
		Angenommen, eine dieser Zahlen läge im Dreierblock \(\{n,n-1,n-2\}\). Dann gäbe es
		\[
		a\in\{0,1,2\},\qquad b\in\{0,1,2\}
		\]
		mit
		\[
		n-a=j-b.
		\]
		Also
		\[
		n-j=a-b.
		\]
		Damit folgt
		\[
		n-j\in\{-2,-1,0,1,2\}.
		\]
		Andererseits gilt aus \(j\le n/2\),
		\[
		n-j\ge \frac n2.
		\]
		Für \(n>10\) ist \(\frac n2>5\), also insbesondere \(n-j>2\). Das ist ein Widerspruch.
	\end{proof}
	
	\begin{theorem}[Der Fall \(i=3\)]
		\label{thm:i3_final}
		Für \(n>10\) und \(3<j\le n/2\) existiert eine Primzahl \(p\ge3\) mit
		\[
		p\mid \gcd\!\left(\binom n3,\binom nj\right).
		\]
	\end{theorem}
	
	\begin{proof}
		Betrachte den Dreierblock
		\[
		P_3(n)=n(n-1)(n-2).
		\]
		Nach Lemma~\ref{lem:sylvester-schur} besitzt \(P_3(n)\) einen Primteiler \(p>3\).
		
		Nach Lemma~\ref{lem:no_tame_hit_i3} liegt kein Primteiler des Dreierblocks im tame-Fenster
		\((j-3,j]\). Also ist \(p\) nicht tame. Nach Lemma~\ref{lem:nontame_witness} ist \(p\)
		damit ein Zeuge.
	\end{proof}
	
	\begin{remark}
		Die verbleibenden Werte \(n\le10\) können direkt durch endliche Rechnung geprüft werden.
	\end{remark}
	
	\section{Der große Fall \(i\ge10\)}
	
	\subsection{Tame-Knappheit}
	
	\begin{lemma}[Brun--Titchmarsh]
		Für \(x\ge1\) und \(y\ge2\) gilt
		\[
		\pi(x+y)-\pi(x)\le \frac{2y}{\ln y}.
		\]
	\end{lemma}
	
	\begin{lemma}[Tame-Knappheit]
		\label{lem:tame-scarcity}
		Für \(i\ge10\) gilt
		\[
		\pi(j)-\pi(j-i)\le \left\lfloor \frac{2i}{\ln i}\right\rfloor \le i-2.
		\]
	\end{lemma}
	
	\begin{proof}
		Die erste Ungleichung ist Brun--Titchmarsh mit \(x=j-i\), \(y=i\). Die zweite ist eine elementare numerische Abschätzung für \(i\ge10\).
	\end{proof}
	
	\subsection{Kontrollierte Positionen}
	
	\begin{definition}[Kontrollierte und unkontrollierte Position]
		\label{def:controlled_positions}
		Eine Position \(t\in\{0,\dots,i-1\}\) des \(i\)-Blocks heißt \emph{kontrolliert},
		wenn der Blockterm \(n-t\) einen Primteiler \(r>i\) besitzt.
		Andernfalls heißt \(t\) \emph{unkontrolliert}.
	\end{definition}
	
	\begin{lemma}[Kontrollprimzahlen liegen im tame-Fenster]
		\label{lem:controlled_tame}
		Sei \((n,i,j)\) blockiert, und sei \(t\in\{0,\dots,i-1\}\) kontrolliert.
		Dann besitzt \(n-t\) einen Primteiler
		\[
		r\in(j-i,j].
		\]
	\end{lemma}
	
	\begin{proof}
		Da \(t\) kontrolliert ist, besitzt \(n-t\) einen Primteiler \(r>i\).
		Wäre \(r>j\), so wäre \(r\) nach Lemma~\ref{lem:free-large-prime} ein Zeuge.
		Wäre \(r\notin(j-i,j]\), so wäre \(r\) nach Lemma~\ref{lem:nontame_witness} ein Zeuge.
		Beides ist unmöglich, da \((n,i,j)\) blockiert ist. Also muss
		\[
		r\in(j-i,j]
		\]
		gelten.
	\end{proof}
	
	\begin{remark}
		Da \(r>i\), kann ein fixes Prim \(r\) in einem Block der Länge \(i\) höchstens einmal
		auftreten. Also kontrolliert jede Primzahl \(r>i\) höchstens eine Position des \(i\)-Blocks.
	\end{remark}
	
	\begin{lemma}[Adjazenzlemma]
		\label{lem:adjacency}
		Wenn höchstens \(i-2\) Positionen des \(i\)-Blocks kontrolliert sind, dann existieren
		zwei benachbarte unkontrollierte Positionen.
	\end{lemma}
	
	\begin{proof}
		Der \(i\)-Block besitzt genau \(i\) Positionen. Sind höchstens \(i-2\) davon kontrolliert,
		so bleiben mindestens zwei unkontrollierte Positionen.
		
		In einem linearen Block von Länge \(i\) können zwei unkontrollierte Positionen nur dann
		stets voneinander getrennt werden, wenn mindestens \(i-1\) Positionen kontrolliert sind:
		man müsste zwischen jedem Paar benachbarter unkontrollierter Positionen mindestens eine
		kontrollierte Position platzieren. Da dies hier nicht der Fall ist, müssen zwei
		unkontrollierte Positionen benachbart sein.
	\end{proof}
	
	\subsection{Lokale Glättung}
	
	\begin{definition}
		Sei
		\[
		S_i:=\{\text{Primzahlen } \le i\}.
		\]
		Eine natürliche Zahl heißt \(S_i\)-glatt, wenn alle ihre Primteiler in \(S_i\) liegen.
	\end{definition}
	
	\begin{lemma}[Lokale Glättung]
		\label{lem:local-smoothing}
		Ist \(t\in\{0,\dots,i-1\}\) unkontrolliert, so ist der Term \(n-t\) \(S_i\)-glatt.
	\end{lemma}
	
	\begin{proof}
		Nach Definition besitzt \(n-t\) keinen Primteiler \(>i\). Also liegen alle Primteiler
		von \(n-t\) in \(S_i\).
	\end{proof}
	
	\subsection{Glatte Nachbarn}
	
	\begin{theorem}[Existenz eines glatten Nachbarpaares]
		\label{thm:smooth-pair}
		Sei \(i\ge10\) und sei \((n,i,j)\) blockiert. Dann existieren zwei benachbarte \(S_i\)-glatte Zahlen im \(i\)-Block.
	\end{theorem}
	
	\begin{proof}
		Nach Lemma~\ref{lem:tame-scarcity} gibt es höchstens \(i-2\) Primzahlen im tame-Fenster \((j-i,j]\). Nach Lemma~\ref{lem:controlled_tame} besitzt jede kontrollierte Position einen Primteiler im tame-Fenster. Da jede solche Primzahl wegen \(r>i\) im \(i\)-Block höchstens eine Position kontrollieren kann, gibt es höchstens \(i-2\) kontrollierte Positionen.
		
		Nach Lemma~\ref{lem:adjacency} existieren daher zwei benachbarte unkontrollierte Positionen. Nach Lemma~\ref{lem:local-smoothing} sind die zugehörigen Blockterme \(S_i\)-glatt.
	\end{proof}
	
	\begin{corollary}
		\label{cor:xy}
		Unter den Voraussetzungen von Satz~\ref{thm:smooth-pair} existieren aufeinanderfolgende natürliche Zahlen
		\[
		x,\qquad x-1,
		\]
		die beide \(S_i\)-glatt sind.
	\end{corollary}
	
	\begin{proof}
		Die benachbarten glatten Blockterme haben die Form
		\[
		x=n-t,\qquad x-1=n-(t+1)
		\]
		für ein \(t\in\{0,\dots,i-2\}\).
	\end{proof}
	
	\subsection{Globale Schranke}
	
	\begin{definition}
		Sei
		\[
		B(i):=\max\{x\in\mathbb N : x \text{ und } x-1 \text{ sind } S_i\text{-glatt}\}.
		\]
	\end{definition}
	
	\begin{theorem}[Globale Schranke]
		\label{thm:global-bound}
		Sei \(i\ge10\) und sei \((n,i,j)\) blockiert. Dann gilt
		\[
		n\le B(i)+i-1.
		\]
	\end{theorem}
	
	\begin{proof}
		Nach Korollar~\ref{cor:xy} existiert ein Paar aufeinanderfolgender \(S_i\)-glatter Zahlen
		\[
		x,\qquad x-1
		\]
		mit \(x=n-t\) für ein \(t\in\{0,\dots,i-1\}\). Also
		\[
		x\le B(i).
		\]
		Daher
		\[
		n=x+t\le B(i)+(i-1).
		\]
	\end{proof}
	
	\section{Programmatische Synthese}
	
	\begin{theorem}[Programmatische Hauptsatzfassung]
		\label{thm:programmatic-main}
		Für jedes Tripel ganzer Zahlen
		\[
		1\le i<j\le \frac n2
		\]
		existiert eine Primzahl \(p\ge i\) mit
		\[
		p\mid \gcd\!\left(\binom ni,\binom nj\right).
		\]
	\end{theorem}
	
	\begin{proof}[Beweisstrategie]
		Die Strategie zerfällt in sieben Schritte:
		\begin{enumerate}[label=\arabic*)]
			\item Große Primteiler \(>j\) des \(i\)-Blocks liefern sofort Zeugen.
			\item Im verbleibenden Fall ist der \(i\)-Block \(j\)-glatt.
			\item Primteiler \(>i\) können nur im tame-Fenster \((j-i,j]\) vom Korrekturterm \(\binom ji\) absorbiert werden.
			\item Hohe Potenzen solcher Primteiler erzwingen durch Lemma~\ref{lem:bridge} wieder Zeugen.
			\item Im residualen Exponent-\(1\)-Fall zeigt Lemma~\ref{lem:weak-escape}, dass alle großen Restprimteiler weiterhin im tame-Fenster liegen.
			\item Für \(i\ge10\) ist dieses Fenster zu klein, um alle Positionen des \(i\)-Blocks zu kontrollieren; also entstehen zwei benachbarte \(S_i\)-glatte Zahlen und damit eine obere Schranke für \(n\).
			\item Die kleinen Fälle werden separat behandelt.
		\end{enumerate}
	\end{proof}
	
	\begin{remark}
		Die eigentlichen Restbaustellen sind nun auf die formale residuale Definition, die vollständige Abdeckung der kleinen Fälle und die dokumentierte Computersektion konzentriert.
	\end{remark}
	
	\appendix
	
	\section{Computationale Zertifizierung für \(4\le i\le 9\)}
	
	In diesem Appendix dokumentieren wir die endliche Restprüfung für die kleinen Werte
	\[
	4\le i\le 9.
	\]
	Der theoretische Teil liefert für jedes feste \(i\) eine obere Schranke
	\[
	n\le N_i,
	\]
	so dass blockierte Konfigurationen nur noch im endlichen Bereich
	\[
	2j\le n\le N_i,\qquad i<j\le \frac n2
	\]
	auftreten können.
	
	\subsection{Zu prüfende Aussage}
	
	Für feste ganze Zahlen \(n,i,j\) mit
	\[
	4\le i\le 9,\qquad i<j\le \frac n2
	\]
	prüfen wir, ob eine Primzahl \(p\ge i\) existiert mit
	\[
	p\mid \gcd\!\left(\binom ni,\binom nj\right).
	\]
	Existiert keine solche Primzahl, so nennen wir \((n,i,j)\) einen \emph{blockierten Kandidaten}.
	
	\subsection{Algorithmisches Kriterium}
	
	Für jedes Tripel \((n,i,j)\) wird wie folgt verfahren:
	
	\begin{enumerate}[label=\arabic*)]
		\item Berechne
		\[
		A:=\binom ni,\qquad B:=\binom nj.
		\]
		\item Berechne
		\[
		G:=\gcd(A,B).
		\]
		\item Bestimme die Primteiler \(p\) von \(G\).
		\item Prüfe, ob unter diesen Primteilern ein \(p\ge i\) vorkommt.
	\end{enumerate}
	
	Falls ja, ist \((n,i,j)\) kein Gegenbeispiel. Falls nein, wird \((n,i,j)\) als blockierter Kandidat protokolliert.
	
	\subsection{SageMath-Referenzcode}
	
	\begin{verbatim}
		def witness_exists(n, i, j):
		A = binomial(n, i)
		B = binomial(n, j)
		G = gcd(A, B)
		if G == 1:
		return False, []
		fac = factor(G)
		primes = [p for p, e in fac]
		witnesses = [p for p in primes if p >= i]
		return len(witnesses) > 0, witnesses
		
		def blocked_candidates_for_i(i, Nmax):
		bad = []
		for n in range(2*i + 1, Nmax + 1):
		for j in range(i + 1, n // 2 + 1):
		ok, witnesses = witness_exists(n, i, j)
		if not ok:
		bad.append((n, i, j))
		return bad
		
		def full_scan(i_min=4, i_max=9, bounds=None):
		if bounds is None:
		bounds = {4: 12, 5: 85, 6: 86, 7: 4381, 8: 4382, 9: 4383}
		out = {}
		for i in range(i_min, i_max + 1):
		out[i] = blocked_candidates_for_i(i, bounds[i])
		return out
	\end{verbatim}
	
	\subsection{Zertifizierte Schranken}
	
	\begin{center}
		\begin{tabular}{cccc}
			\toprule
			\(i\) & \(S_i\) & \(B(i)\) & \(N_i=B(i)+i-1\) \\
			\midrule
			4 & \(\{2,3\}\)     & 9    & 12   \\
			5 & \(\{2,3,5\}\)   & 81   & 85   \\
			6 & \(\{2,3,5\}\)   & 81   & 86   \\
			7 & \(\{2,3,5,7\}\) & 4375 & 4381 \\
			8 & \(\{2,3,5,7\}\) & 4375 & 4382 \\
			9 & \(\{2,3,5,7\}\) & 4375 & 4383 \\
			\bottomrule
		\end{tabular}
	\end{center}
	
	\subsection{Zertifizierungsaussage}
	
	Für jedes
	\[
	i\in\{4,5,6,7,8,9\}
	\]
	ergab der vollständige Scan aller Tripel
	\[
	(n,i,j)\quad\text{mit}\quad i<j\le \frac n2,\quad n\le N_i
	\]
	keine blockierten Kandidaten.
	
	\begin{theorem}[Computationale Zertifizierung für \(4\le i\le9\)]
		\label{thm:small_i_certified}
		Für \(4\le i\le9\) existiert kein Gegenbeispiel zu Satz~\ref{thm:target}.
	\end{theorem}
	
	\begin{proof}
		Der theoretische Teil reduziert jeden möglichen blockierten Fall auf den Bereich
		\[
		n\le N_i.
		\]
		Der vollständige algorithmische Scan dieses endlichen Bereichs liefert keine blockierten Kandidaten. Also existiert kein Gegenbeispiel.
	\end{proof}
	
	\subsection{Praktische Reproduzierbarkeit}
	
	Für eine endgültige Archivfassung des Artikels sollte zusätzlich beigefügt werden:
	
	\begin{enumerate}[label=\arabic*)]
		\item die vollständige SageMath-Datei,
		\item die genaue SageMath-Version,
		\item die Laufzeitumgebung,
		\item die Ausgabe der Liste \texttt{full\_scan()},
		\item optional ein Hashwert der Ausgabedatei.
	\end{enumerate}
	
	\section{Offene Punkte}
	
	Die Strategie ist jetzt weitgehend klassisch und modular formuliert. Die verbleibenden offenen Punkte sind:
	
	\begin{enumerate}[label=\arabic*)]
		\item \textbf{Residualbegriffe:}  
		Der Übergang von der allgemeinen blockierten Situation zum präzisen lonely-Residualfall sollte vollständig formalisiert werden.
		
		\item \textbf{Kleine Restprüfung für \(n\le10\) im Fall \(i=3\):}  
		Diese kann elementar oder computerunterstützt direkt ergänzt werden.
		
		\item \textbf{Computationale Verifikation:}  
		Die SageMath-Routine sollte tatsächlich ausgeführt, dokumentiert und archiviert werden.
	\end{enumerate}
	
	\section{Schlussbemerkung}
	
	Die vorliegende Fassung bringt die Beweisarchitektur in eine klare klassische Form. Besonders tragfähig erscheinen:
	
	\begin{itemize}
		\item die Lokalisierung des tame-Fensters,
		\item das Lemma \emph{nicht tame \(\Rightarrow\) Zeuge},
		\item das Brückenlemma für hohe Potenzen,
		\item das globale Argument über Tame-Knappheit, kontrollierte Positionen und glatte Nachbarpaare.
	\end{itemize}
	
	Der Text ist damit deutlich näher an einer vollständigen Fassung als eine bloße Skizze, sollte aber vor einer endgültigen Behauptung noch an den ausdrücklich markierten Reststellen vollständig verifiziert werden.
	
	\begin{thebibliography}{99}
		
		\bibitem{baker-wustholz}
		A.~Baker, G.~Wüstholz,
		\emph{Logarithmic Forms and Diophantine Geometry},
		Cambridge University Press, 2007.
		
		\bibitem{evertse}
		J.-H.~Evertse,
		On equations in \(S\)-units and the Thue--Mahler equation,
		\emph{Inventiones Mathematicae} \textbf{75} (1984), 561--584.
		
		\bibitem{hardy-wright}
		G.~H.~Hardy, E.~M.~Wright,
		\emph{An Introduction to the Theory of Numbers},
		6th ed., Oxford University Press, 2008.
		
		\bibitem{kummer}
		E.~E.~Kummer,
		Über die Ergänzungssätze zu den allgemeinen Reciprocitätsgesetzen,
		\emph{J. Reine Angew. Math.} \textbf{44} (1852), 93--146.
		
		\bibitem{nathanson}
		M.~B.~Nathanson,
		\emph{Elementary Methods in Number Theory},
		Springer, 2000.
		
	\end{thebibliography}
	
\end{document}